\section{The record: from case files to dated events}\label{sec:record}

\subsection{Corpus and parse}

The USPTO Trademark Annual Applications backfile (product TRTYRAP) is the
complete record of US trademark filings from 1884 onward: 13.99 million
case files, each carrying the goods/services description by class, the
prosecution history, and the ownership and correspondence records. The
prosecution history is a list of coded events with dates. We re-parse
every case file into 242 million dated events, keyed by serial number,
so that each application's full sequence is available: office actions,
responses, allowances, statements of use, registration, Section~8 and
Section~9 filings, cancellations and renewals, each with its
date.\footnote{Event codes are read against the USPTO's own event-code
dictionary. Terminal status codes are retained for cross-checking: 700
registered; 701--705 Section~8 accepted; 710 cancelled for Section~8
failure; 800 renewed; 900 expired. The year-six cancellation and the
year-ten cancellation share status 710 and are separable only by event
date, which is the reason the corpus is event-dated.} The parse also
reads counsel of record on each filing, the statutory filing basis (use,
intent to use, \S44(e) home registration, \S66(a) Madrid extension), the
owner's domicile, and the owner name, which is normalized ---
uppercased, stripped of punctuation and of two dozen legal suffixes ---
so that \textsc{acme corp.} and \textsc{acme corporation} are one owner.
\citet{Graham2013} describe the administrative data; the event-dated
parse adds the full prosecution timeline of each case.

\subsection{The steps, the samples, and the conventions}\label{sec:gates}

A US trademark passes through a sequence of dated, high-stakes steps.
An application is filed with a goods/services description. An examiner
reviews it, often issuing office actions that the applicant must answer.
If the examiner allows the application --- and, for intent-to-use
filings, once the applicant files a sworn statement that the mark has
entered commerce --- the mark registers. Registration then starts a
maintenance clock. In years five to six the owner must file a Section~8
declaration swearing that the mark is still in use in commerce, with a
six-month grace period; this is the \emph{five-year proof of continued
use}, the paper's central outcome. Around year ten a combined Section~8
declaration and Section~9 renewal is due, and again every ten years
after that. A mark whose owner misses a proof is cancelled. Madrid
Protocol registrations (\S66(a) basis) file Section~71 declarations on
the same schedule and are flagged separately throughout.

\begin{table}[h]\centering\small
\caption{The trademark funnel in this corpus. Registration counts are
unique serials; five-year-proof outcomes are event-dated (a Section~8
cancellation at registration age 4.0--8.5 years). Source: the re-parsed
USPTO application backfile.}
\label{tab:funnel}
\begin{tabular}{lrr}
\toprule
Stage & $n$ & Share \\
\midrule
Debut filings scored, 1995--2018 (owners' first filings) & 3{,}589{,}068 & \\
\quad reached registration & 2{,}066{,}600 & 57.6\% \\
\midrule
All unregistered class-records filed 1995--2018 & 3{,}798{,}318 & \\
\quad never filed statement of use & & 42.6\% \\
\quad stopped responding in prosecution & & 45.3\% \\
\quad removed adversarially & & 4.3\% \\
\midrule
Registrations 2002--2018, scored, unique serials & 3{,}104{,}534 & \\
\quad failed the five-year proof of continued use (event-dated) & & 52.3\% \\
\midrule
Passed the five-year proof, cohorts 2002--2013, status known & 1{,}129{,}724 & \\
\quad failed at the year-ten renewal & & 39.1\% \\
\bottomrule
\end{tabular}
\end{table}

Three named outcomes recur (Table~\ref{tab:funnel}). \emph{Failing to
register}: the application never completed, which happens by the
applicant's own inaction in roughly nine cases of ten. \emph{Failing the
five-year proof of continued use}: a registration cancelled for non-use
at registration age 4.0 to 8.5 years, which means the product the mark
named left commerce within six years. \emph{Reaching the SEC record}:
the owner appears in dated SEC events.

Outcomes are rates, so every estimate is a shift in a probability and is
stated in percentage points against the base rate it shifts. The basic
comparison is the outcome rate of the fifth of filings with the highest
lead minus that of the fifth with the lowest, with the fifths cut inside
a single Nice class and year so that no difference between industries or
between years enters; regression tables add drafting and filer controls
and cluster errors on the owner. A filing may claim several Nice classes
at once, so a \emph{class-record} counts it once per class while a
\emph{registration} counts it once by serial number; a \emph{debut
filing} is an owner's first; a \emph{cohort} is the set of registrations
issued in a calendar year.

Two windows bound the samples, each set by an institutional fact.
Scoring runs 1995--2019, because a filing can only be scored when the
record holds a full five years of its class's language on either side of
its date. Full windows exist for filings made 1990 through 2020; scoring
holds back five years at the start, where references are thinnest, and
one at the end. The supplementary
material (\S S3) shows how thin references distort scores near the
start of the record. Five-year-proof
cohorts run 2002--2018: cancellations post at about registration age six
and a half, and the record ends in April 2026, so 2018 is the last
cohort old enough for its proof window to have essentially elapsed.

\subsection{Owner links}\label{sec:links}

Owners are matched to three external records by normalized exact name.
\emph{SEC reporting} is any appearance of the owner as a registrant in
the Commission's EDGAR filer index; roughly half the matched firms carry
a ticker and the rest report for other reasons. \emph{Regulation~D}
rounds are Form~D notices of exempt offerings, electronic and therefore
complete only from March 2009. \emph{Patents} are PatentsView assignee
records \citep{PatentsView}. The SEC match runs over all 5.67 million
owners in the corpus; names resolving to more than one registrant are
dropped rather than assigned, yielding 19{,}889 matched owners covering
7{,}588 registrants. Validation on 22 companies that indisputably went
public matches all 22, and a hand inspection of 25 random matches found
none wrong --- a check on gross failure modes rather than a precision
estimate. The linkage is nonetheless a name match, and that bounds what
can be read from it: distinctive names are easier to match, and
distinctiveness of a name plausibly travels with distinctiveness of the
text. Section~\ref{sec:validation} returns to which results survive
that bound.
